\par\Needspace{6\baselineskip}\medskip
\hypertarget{native-217-article158}{}\pdfbookmark[2]{158.}{article-158}
\begin{center}158.\end{center}
Die beiden vorhergehenden Beispiele gaben noch keinen Anlass, die Methode aus Artikel~120 anzuwenden: Die aufeinanderfolgenden Hypothesen konvergierten so rasch, dass man schon bei der zweiten hätte stehen bleiben können und die dritte kaum merklich vom wahren Ergebnis abwich. Diesen Vorteil haben wir tatsächlich immer und können auf eine vierte Hypothese verzichten, wenn die heliozentrische Bewegung mäßig ist und die drei Radiusvektoren nicht allzu verschieden sind, besonders wenn außerdem die Zeitintervalle wenig voneinander abweichen. Je weiter sich die Bedingungen der Aufgabe jedoch davon entfernen, desto stärker weichen die zuerst angenommenen Werte der Größen $P$, $Q$ von den wahren ab, und desto langsamer konvergieren die folgenden Werte gegen die wahren. In einem solchen Fall sind die ersten drei Hypothesen so durchzurechnen, wie es die beiden vorhergehenden Beispiele zeigen (mit dem einzigen Unterschied, dass man in der dritten Hypothese wie in der ersten und zweiten die
\SourceLeaf{218}
Größen $\eta$, $\eta''$, $P'$, $Q'$, $X$, $Y$ anstelle der Elemente selbst berechnen muss). Dann werden die zuletzt erhaltenen Werte von $P'$, $Q'$ nicht mehr als neue Werte der Größen $P$, $Q$ in der vierten Hypothese angesetzt. Diese neuen Werte ermittelt man vielmehr nach der Methode aus Artikel~120 durch Kombination der ersten drei Hypothesen. Nur sehr selten wird es dann nötig sein, nach den Anweisungen in Artikel~121 zu einer fünften Hypothese überzugehen. --- Wir veranschaulichen nun auch diese Rechnungen an einem Beispiel, das zugleich zeigen wird, wie weit unsere Methode anwendbar ist.
